\begin{figure}[H]
\centering
\includegraphics[width=1.0\linewidth]{wakayama/M1-200_month_grid_200hPa.pdf}
\caption[200 hPa の月別分布図]{\SI{200}{\hecto\pascal}（約 \SI{12}{\kilo\metre}）の平均風を月ごとに並べたもの。塗りが風速、矢印が風向。ジェットが季節とともに動く様子が地図の上で見える。}
\label{fig:m1200-wakayama}
\end{figure}

\begin{figure}[H]
\centering
\includegraphics[width=1.0\linewidth]{wakayama/M1-030_month_grid_30hPa.pdf}
\caption[30 hPa の月別分布図]{同じく \SI{30}{\hecto\pascal}（約 \SI{24}{\kilo\metre}）。対流圏とは逆の季節変化をしており、夏に東風が強まる。}
\label{fig:m1030-wakayama}
\end{figure}

\begin{figure}[H]
\centering
\includegraphics[width=1.0\linewidth]{wakayama/M3-200_direction_map_200hPa.pdf}
\caption[風向の分布図]{\SI{200}{\hecto\pascal} の平均風向。円環用のカラーマップを使い、定常度 0.3 未満は灰色にしている。灰色の領域では平均風向に意味がない。}
\label{fig:m3200-wakayama}
\end{figure}

\begin{figure}[H]
\centering
\includegraphics[width=1.0\linewidth]{wakayama/M2_constancy_map.pdf}
\caption[定常度の分布図]{高度と季節ごとの定常度。濃いところでは風向が安定していて計画が立てやすく、淡いところではビンのなかで風向が入れ替わる。}
\label{fig:m2-wakayama}
\end{figure}

\begin{figure}[H]
\centering
\includegraphics[width=0.92\linewidth]{wakayama/M4-01_wind_ellipse_map.pdf}
\caption[風のばらつきを地図に]{1 月前半、\SI{200}{\hecto\pascal} の平均風ベクトルと、その先端に描いた $1\sigma$ の散らばり楕円。丸い楕円はどの向きにもばらつくこと、細長い楕円は主に強さがばらつくことを意味する。}
\label{fig:m401-wakayama}
\end{figure}
